\section{Resuming selected payments without losing their history}
\label{sim:coupon-program-resume}

A selected same-period runner now extends the checked coupon prefix through several supplied instructions. Its plan retains the original period opening for reports and the current book for execution. It also identifies the original source, selected method registry, verified prefix, complete future instruction list and resource limits. A checkpoint's cursor counts only accepted new instructions. The earlier coupon refusal remains in its archive; it is not replaced by the outcome of a later payment.

For a successful three-instruction suffix, processing one event per call, two per call or the whole suffix produces the same terminal checkpoint. The complete accepted journals, dated flows, current book and cumulative report agree with independent arithmetic and uninterrupted execution. This comparison also covers a later unaffordable payment, duplicate identity and backwards date. In each case, execution retains the accepted prefix and leaves the refused instruction first in the remaining list.\footnote{\path{data/accounting-master/coupon-program-resume-20261027/} retains sixteen cases and 48 chunk comparisons, with 824 assertions and eighteen adverse controls. A separate audit adds 493 assertions, twelve controls and two exact byte-boundary comparisons. The first attempt is preserved: its comparator incorrectly required identical duplicate-refusal wrappers. The generic runner detects a duplicate command before calling the event engine; the selected runner retains the engine's duplicate-event diagnostic. Both reject the same instruction and preserve the same financial state.}

An execution pause has a different meaning from a refusal. A call may consume its allowed number of instructions while another valid instruction remains; this returns a budget pause and can continue from the checked cursor. An economic or chronology refusal stops before the offending event, which cannot be silently skipped on a subsequent call. Changing that event requires a newly identified plan. A complete, refused or resource-limited checkpoint remains unchanged when submitted again under its existing plan.

Resource limits apply to the entire selected suffix, not afresh to each chunk. The inherited command ceiling is 64 and the default trace allowance is 8 MiB. A deliberately smaller allowance exposes the boundary exactly: 82,883 bytes fit the first accepted trace row, but not the second. An allowance of 82,882 bytes accepts no row. The runner retains the next command and reports the unchanged accepted state in both cases. Splitting execution into additional calls does not enlarge the allowance. These limits govern this diagnostic implementation; they neither determine the financial shortfall nor prove that a larger economically feasible schedule is impossible.

The checkpoint's hash is necessary for identifying the supplied record but is insufficient for validating it. The selected interface recomputes the bounded accepted prefix from the identified source, registry and instruction list, then compares the complete checkpoint. This detects altered request metadata even when a narrower financial comparison would still find the same balances. It also detects a changed opening, skipped cursor, missing trace, erased refusal, modified report or enlarged plan budget. Recomputed validation does not append a second accepted event or duplicate a cash flow.

This evidence establishes a selected same-period partial-coupon interface. It does not migrate the general program runner, support every method or reporting-period transition, or establish durable exactly-once commits under concurrent writers. Repeated prefix computation is also a cost, not a scalability result. The distinction between original opening, accepted current book and remaining instructions nevertheless supplies a concrete starting point for broader resumption contracts without erasing the company's earlier history.
